\chapter{DFT}

\begin{figure}[h]
    \centering
    \includegraphics[scale = 0.35]{Fourier_transform,_Fourier_series,_DTFT,_DFT.png}
\end{figure}  

\newpage 

\section{Sequenze periodiche}

\footnote{ 
    Slide del prof | 7 - DFT | pag 2 - 3\\
    Appunti| 2025-10-28 | pag 4   
}

Consideriamo una sequenza periodica $\tilde{x} \left[ n\right]$ per cui: 

{
    \Large 
    \begin{equation}
        \tilde{x} \left[ n\right] 
        = 
        \tilde{x} \left[ n + N\right] 
    \end{equation}
}

dove N è un numero intero. \newline 

\begin{tcolorbox}
    A livello di notazione, scriveremo $\tilde{x} \left[ n\right]$, cioè una sequenza $x \left[ n\right]$ periodica, 
    se è presente un'ondina \~ sopra alla sequenza. 
\end{tcolorbox}

Di una sequenza periodica $\tilde{x} \left[ n\right]$ possiamo determinare la sua pulsazione (fondamentale) $\omega_0$ come: 

{
    \Large 
    \begin{equation}
        \omega_0 = \dfrac{2 \pi}{N}
    \end{equation}
}

A titolo di esempio, possiamo utilizzare la seguente figuta: 

\begin{figure}[h]
    \centering
    \includegraphics[scale = 0.3]{Sequenza periodica di esempio.png}
\end{figure} 

Come si nota dalla figura, la conoscenza degli N campioni, cioè: 

{
    \Large 
    \begin{equation}
        \tilde{x} \left[ k\right] \text{ per } k = 0, 1, \dots, \left(N - 1\right)
    \end{equation}
}

permette di conoscere l'intera sequenza periodica. \newline 

Proprio perchè una sequenza periodica $\tilde{x} \left[ n\right]$ non è nè assoltamente sommabile, 
nè quadraticamente sommabile, allora, per calcolare la sua DTFT la calcoleremo come: 

{
    \Large 
    \begin{equation}
        \tilde{x} \left[ n\right]
        = 
        x \left[ n\right]
        *
        s \left[ n\right]
    \end{equation}
}

cioè una convoluzione tra $x \left[ n\right]$ e $s \left[ n\right]$, dove: 

{
    \Large 
    \begin{equation}
        s \left[ n\right] 
        = 
        \sum_{k = - \infty}^{+ \infty}
        \delta \left[ n - k N\right]
    \end{equation}
}

{
    \Large 
    \begin{equation}
        x \left[ n\right]
        = 
        \begin{cases}
            \tilde{x} \left[ n\right] \text{ per } 0 \leq n \leq N - 1 
            \\
            0 \text{ altrove}
        \end{cases}
    \end{equation}
}

\newpage 

In altre parole, una sequenza periodica $\tilde{x} \left[ n\right]$ è una convoluzione di un periodo con un treno di impulsi $s \left[ n\right]$ di periodo N. \newline 

Da un punto di vista grafico, a titolo di esempio: 

\begin{figure}[h]
    \centering
    \includegraphics[scale = 0.25]{DTFT di una sequenza periodica esempio.png}
\end{figure} 

\newpage 

\subsection{DFS - Discrete Fourier series}

\footnote{ 
    Slide del prof | 7 - DFT | pag 4 - 5\\
    Appunti| 2025-10-28 | pag 4   
}

Come nel caso continuo, le sequenze periodiche possono essere rappresentate dalla serie discreta di Fourier (a.k.a. DFS - Discrete Fourier Series). \newline 

Quindi, possiamo rappresentare una qualasiasi sequenza periodica $\tilde{x} \left[ n\right]$ che si estende da $\left[ - \infty, + \infty\right]$ come: 

{
    \Large 
    \begin{equation}
        \tilde{x} \left[ n\right]
        = 
        \dfrac{1}{N}
        \sum_{k}
        \tilde{X} \left[ k\right] \cdot e^{\jmath \omega_0 k n}
    \end{equation}
}

dove: 

\begin{itemize}
    \item N è il periodo 
    \item $\dfrac{1}{N}$ è un coefficiente moltiplicativo 
    \item $\tilde{X} \left[ k\right]$ è il coefficiente della serie, cioè il k-esimo coefficiente della DFS di $\tilde{x} \left[ n\right]$
\end{itemize}

Dalla formula di $\tilde{x} \left[ n\right]$, possiamo svolgere diverse considerazioni. \newline 

Grazie anche alla notazione, dalla formula di $\tilde{x} \left[ n\right]$, possiamo notare che i coefficiente della serie DFS, 
$\tilde{X} \left[ k\right]$ sono estendibili da $\left[ - \infty, + \infty\right]$, proprio come $\tilde{x} \left[ n\right]$. \newline 

Sapendo che la pulsazione fondamentale $\omega_0$ vale: 

{
    \Large 
    \begin{equation}
        \omega_0 
        =
        \dfrac{2 \pi}{N}
    \end{equation}
}

esistono solo N distinti esponenziali complessi: 

{
    \Large 
    \begin{equation}
    \sum_{k = 0}^{N - 1}
    e^{\jmath \omega_0 k n} 
    = 
    e^{\jmath \frac{2 \pi}{N} \cdot 0}
    +
    e^{\jmath \frac{2 \pi}{N} \cdot 1}
    + 
    \dots 
    e^{\jmath \frac{2 \pi}{N} \cdot \left(N  - 1\right)}  
    \end{equation}
}

Grazie alla periodicità di $\tilde{X} \left[ k\right]$ e $\tilde{x} \left[ n\right]$, possiamo scrivere che: 

{
    \Large 
    \begin{equation}
        e^{\jmath \frac{2 \pi}{N} \cdot 0} 
        = 
        e^{\jmath \frac{2 \pi}{N} \cdot N}
    \end{equation}
}

In generale, dato un valore di k qualasiasi, possiamo scrivere che: 

{
    \Large 
    \begin{equation}
        e^{\jmath \frac{2 \pi}{N} \cdot k} 
        = 
        e^{\jmath \frac{2 \pi}{N} \cdot \left(k + N\right)}
    \end{equation}
}

Scrivendo $e^{- \jmath \frac{2 \pi}{N}}$ come: 

{
    \Large 
    \begin{equation}
        W_N = e^{- \jmath \frac{2 \pi}{N}}
    \end{equation}
}

allora possiamo scrivere $\tilde{x} \left[ n\right]$ come: 

{
    \Large 
    \begin{equation}
        \begin{split}
        \tilde{x} \left[ n\right]
        &= 
        \dfrac{1}{N}
        \sum_{k}
        \tilde{X} \left[ k\right] \cdot e^{\jmath \omega_0 k n}
        \\
        &\downarrow
        \\
        \tilde{x} \left[ n\right]
        &= 
        \dfrac{1}{N}
        \sum_{k}
        \tilde{X} \left[ k\right] \cdot W_N^{-k n}
        \end{split}
    \end{equation}
}

\newpage 

Il termine $W_N$ prende il nome di Twiddle. \newline 

\begin{tcolorbox}
    Il fattore di Twiddle verrà apprfondito nei futuri capitoli della FFT, ma se intanto vuoi avere un assaggio dell'importanza del Twiddle: 

    \begin{itemize}
        \item \url{https://www.linkedin.com/pulse/how-fft-algorithm-works-part-5-twiddle-factors-mark-newman/}
        \item \url{https://www.sciencedirect.com/topics/engineering/discrete-fourier-transform}
    \end{itemize}
\end{tcolorbox}

In particolare la formula: 

{
    \Large
    \begin{equation}
        \tilde{x} \left[ n\right]
        =
        \dfrac{1}{N}
        \sum_{k = 0}^{N -1}
        \tilde{X} \left[ k\right] \cdot W_N^{-k n}
    \end{equation}
}

prende il nome di DFS inversa. \newline 

Si può dimostrare anche il processo inverso, cioè la DFS diretta come: 

{
    \Large
    \begin{equation}
        \tilde{X} \left[ k\right]
        =
        \sum_{k = 0}^{N -1}
        \tilde{x} \left[ n\right] \cdot W_N^{k n}
    \end{equation}
}

ricordando che: 

{
    \Large 
    \begin{equation}
        \sum_{k = 0}^{N - 1}
        W_N^{kn} 
        = 
        \begin{cases}
            N \text{ per n multiplo di N}
            \\
            0 \text{ altrimenti}
        \end{cases}
    \end{equation}
}

Grazie a questa proprietà dei twiddle $W_N^{kn}$, possiamo dire che il twiddle forma una base ortognale dello spazio nel dominio di DFS. \newline 

\newpage 

\subsubsection{Relazione tra DTFT e DFS}

\footnote{ 
    Slide del prof | 7 - DFT | pag 7 - 8\\
    Appunti| 2025-10-28 | pag 4   
}

Poichè una sequenza periodica $\tilde{x} \left[ n\right]$ con periodo N si può vedere come la convoluzione del singolo periodo di $x \left[ n\right]$ 
con un treno di impulsi $s \left[ n\right]$, spaziati di N, allora, possiamo scrivere che la DTFT di $x \left[ n\right]$ vale: 

{
    \Large 
    \begin{equation}
        X \left(e^{\jmath \omega}\right) 
        = 
        \sum_{n = - \infty}^{+ \infty}
        x \left[ n\right]
        \cdot 
        e^{- \jmath \omega n}
        =
        \sum_{n = 0}^{N - 1}
        x \left[ n\right]
        \cdot 
        e^{- \jmath \omega n} 
    \end{equation}
}

per cui: 

{
    \Large 
    \begin{equation}
        \tilde{X} \left[ k\right]
        = 
        \left. 
        X \left(e^{\jmath \omega}\right)
        \right|_{\omega = \dfrac{2 \pi}{N}k}
    \end{equation}
}

dove $\tilde{X} \left[ k\right]$ è la DFS di $\tilde{x} \left[ n\right]$. \newline 

Grazie a quest'ultima relazione tra $\tilde{X} \left[ k\right]$ e $X \left(e^{\jmath \omega}\right)$, 
possiamo osservare che la DFS di $\tilde{x} \left[ n\right]$, cioè $\tilde{X} \left[ k\right]$, corrisponde alla DTFT di $x \left[ n\right]$, cioè $X \left(e^{\jmath \omega}\right)$, 
campionata in N frequenze equispazionate: 

{
    \Large 
    \begin{equation}
        0, \dfrac{2 \pi}{N} \cdot 1, \dfrac{2 \pi}{N} \cdot 2, \dots, \dfrac{2 \pi}{N} \cdot \left(N - 1\right) 
    \end{equation}
}

In altri termini, eseguire la DFS corrisponde a valutare la trasformata-z della sequenza $x \left[ n\right]$, 
cioè $X \left(z\right)$, su N punti equispaziati sul cerchio unitario, come si visualizza dalla seguente figura: 

\begin{figure}[h]
    \centering
    \includegraphics[scale = 0.25]{Cerchio unitario discretizzato per la DFS.png}
\end{figure} 

Inoltre, si può calcolare la DTFT della sequenza $\tilde{x} \left[ n\right]$, cioè $\tilde{X} \left(e^{\jmath \omega}\right)$ come: 

{
    \Large 
    \begin{equation}
        \tilde{X} \left(e^{\jmath \omega}\right) 
        = 
        \sum_{k = - \infty}^{+ \infty}
        \dfrac{2 \pi}{N}
        \cdot 
        X \left(e^{\jmath \omega}\right)
        \cdot 
        \overline{\delta}
        \left(
            \omega 
            - 
            \dfrac{2 \pi}{N} k
        \right) 
    \end{equation}
}

dove: 

\begin{itemize}
    \item $X \left(e^{\jmath \omega}\right)$ è la trasformata di Fourier di $x \left[ n\right]$ nel singolo periodo 
    \item $\overline{\delta}
        \left(
            \omega 
            - 
            \dfrac{2 \pi}{N} k
        \right) $ è un treno di impulsi di Dirac 
\end{itemize}

\newpage 

\section{Sequenze aperiodiche di durata finita}

\footnote{ 
    Slide del prof | 7 - DFT | pag 9 - 10\\
    Appunti| 2025-10-28 | pag 4   
}

Data una sequenza periodica $\tilde{x} \left[ n\right]$, essa è completamente rappresentata dalla sua sequenza DFS $\tilde{X} \left[ k\right]$ 
con $k = 0, 1, \dots, N-1$. \newline 

Invecem data una sequenza aperiodica $x \left[ n\right]$ di lunghezza N, 
quest'ultima può essere resa periodica replicandola sull'asse n, proprio come si vede nelle figure di esempio: 

\begin{figure}[h]
    \centering
    \includegraphics[scale = 0.5]{Confronto tra una sequenza aperiodica e periodica di esempio.png}
\end{figure} 

Allora, la DFS di $\tilde{x} \left[ n\right]$ corrisponde alla DTFT di $x \left[ n\right]$ campionata in N frequenze equispaziate, 
quindi, per conoscere la DFS non bisogna utilizzare la DTFT. \newline 

\begin{tcolorbox}
    Ma tanto, ingegneristicamente parlando, non la utilizzeremo mai la DFS. 
\end{tcolorbox}

Quindi, possiamo scrivere che: 

{
    \Large 
    \begin{equation}
        x \left[ n\right]  
        \leftrightarrow 
        \left.
        X \left(e^{\jmath \omega}\right)
        \right|_{\omega = \frac{2 \pi}{N} k}
    \end{equation}
}

per $n = 0, 1, \dots, N-1$ e $k = 0, 1, \dots, N-1$

\newpage 

\section{Aliasing nel tempo}

\footnote{ 
    Slide del prof | 7 - DFT | pag 13 - 14\\
    Appunti| 2025-10-28 | pag 4 - 5
}

Campionare la DTFT della sequenza $x \left[ n\right]$, cioè $X \left( e^{\jmath \omega}\right)$, in N punti equispaziati sul cerchio unitario, 
equivale ad avere la DFS di una sequenza periodica di periodo N, ottenuta facendo la convoluzione di $x \left[  n \right]$ con un treno di impulsi spaziati di N. \newline 

\begin{tcolorbox}
    La dimostrazione di questo è svolto nelle pagine 11-12 del PDF "7 - DFT", ma essendo le slides facoltative e non di interesse al corso, perchè tanto la DFS non la useremo mai, l'ho salata. 
\end{tcolorbox}

Per un prefissato valore di N e data la lunghezza di $x \left[ n\right]$, possiamo considerare due casi interessanti: \newline 

\begin{itemize}
    
\item \textbf{$x \left[ n\right]$ di durata finita minore o uguale ad N} \newline 

Come si può visualizzare dalle seguenti figure di esempio: 

\begin{figure}[h]
    \centering
    \includegraphics[scale = 0.35]{No aliasing per la DFS di esempio.png}
\end{figure} 

$x \left[ n\right]$ può essere calcolato dai N campioni della sua $X \left(e^{\jmath \omega}\right)$, in particolare: 

{
    \Large 
    \begin{equation}
        x \left[ n \right]
        = 
        \begin{cases}
            \tilde{x} \left[ n\right] \text{ per } 0 \leq n \leq N -1 
            \\
            0 \text{ altrove}
        \end{cases}
    \end{equation}
}

\item \textbf{$x \left[ n\right]$ di durata superiore ad  N (anche infinita)} \newline 

Come si può visualizzare dalle seguenti figure di esempio: 

\begin{figure}[h]
    \centering
    \includegraphics[scale = 0.35]{Aliasing per la DFS di esempio.png}
\end{figure} 

$x \left[ n\right]$ non può essere calcolato dai N campioni della sua $X \left(e^{\jmath \omega}\right)$: 
questo fenomeno prende il nome di aliasing temporale. \newline 

\end{itemize}

\newpage 

\section{DFT - Discrete Fourier Transform}

\footnote{ 
    Slide del prof | 7 - DFT | pag 15 - 16\\
    Appunti| 2025-10-29 | pag 1
}
